\section{讨论}
\label{sec:discussion}

\subsection{部署方式与扩展能力}
\label{subsec:deployment-extensibility}

AgentGate的部署位置由工具调用链路中的信任边界决定，而不依赖某一种特定的智能体框架。系统采用统一安全内核，并支持两种互补的接入方式。第一种方式是在智能体宿主进程内通过框架适配器接入，在工具注册、调用执行前和结果写回上下文前触发网关分析。这种方式能够获得可信任务标识、智能体身份与工具返回等较完整的上下文，适合Function Tool以及LangGraph、OpenAI Agents SDK等框架原生工具。第二种方式是将网关作为Sidecar或MCP透明代理部署在智能体宿主与工具服务之间，对多个智能体发出的调用实施统一策略。两种方式将工具声明、调用参数、返回结果和会话标识转换为相同的安全事件，复用事实状态、来源跟踪和状态化检测；普通Agent请求不能自行声明可信授权。

这种设计将系统能力与具体协议解耦。接入新的智能体框架时，适配器只需将工具注册、调用前和返回后三类事件转换为网关对象，不需要修改核心分析逻辑；接入新的工具领域时，则可通过扩展资源目录、工具画像、效果推断器与策略约束描述领域特有的资源和副作用。语义分析器的职责被限制为抽取动作、资源、范围、效果、目的地及外部指令等受限事实，最终决策仍由任务授权上界、确定性约束和会话状态共同产生，因此可通过OpenAI兼容接口替换不同能力与成本的大语言模型，而无需改变策略执行流程。上述扩展路径使AgentGate能够在框架、工具类型和语义模型变化时保持一致的安全表示与判定接口。

当前原型已实现FastAPI Sidecar、Function Tool适配器以及LangGraph和OpenAI Agents SDK的轻量封装，可以验证进程内接入与独立服务接入两种基本形态。对于MCP，现有适配器能够将\texttt{list\_tools}结果转换为内部工具对象，并生成\texttt{tools/call}请求，但尚未覆盖客户端会话管理、远程连接、结果转发与异常重连，因而目前体现的是协议适配能力，而非完整的透明MCP代理。完成这些协议生命周期功能不会改变三个安全模块的输入输出关系，但仍是将原型扩展为通用工具网关所需的系统工作。

\subsection{局限性}
\label{subsec:limitations}

AgentGate已经通过可信任务授权和参数绑定后的效果推断缩小任务边界，但自然语言任务通常只描述最终目标，并不会穷举完成任务所需的全部辅助操作。若授权编译过于严格，身份核验、信息检索和格式转换等正常前置步骤可能被误判为范围扩张；若为兼容这些步骤而放宽权限，智能体又可能借助自行生成的子任务扩大攻击面。AgentDojo-Traj离线实验中，当前确定性编译器将漏报率降至6.82\%，但误阻298个正常步骤，误报率达到34.33\%。这说明编译器虽以外部entitlement限制授权扩张，仍难以仅凭初始任务文本同时获得完整流程知识与精确最小权限。

系统通过控制指令模式、来源标记和参数数据依赖检测外部内容的操控意图。然而，检测到外部指令并不等价于证明后续调用由该指令触发：工具结果可能同时包含正常业务数据和隐式重定向，智能体也可能在若干正常推理步骤之后才执行受影响的操作。InjecAgent的base设置没有统一覆盖短语，当前内容检测对直接危害和数据窃取的召回率分别只有47.25\%与32.54\%。若将所有命令式文本直接视为攻击会显著增加开放内容上的误报，而要求显式数据依赖又会遗漏语义改写、间接引用和延迟触发。因此，后续语义分析器应只抽取受限事实，并由本地证据融合策略决策，而不能直接替代参考监控器。

对于跨调用风险，AgentGate维护字段级敏感对象、父对象来源关系、窗口计数和审批消费状态，能够处理直接传递、文本嵌入、URL/Base64编码、凭证使用与审批重放。但真实调用链中的数据还可能经历摘要、翻译、加密、拆分和重新生成，工具服务也可能产生网关不可观测的内部副作用，使基于已知签名和工具效果的标签传播出现遗漏。当前机制尚未形成覆盖任意变换的通用信息流系统；内存后端也不能直接保证多进程一致性，Redis后端仍依赖部署者正确配置。协议原生重放只验证实现不变量，两个公开基准也只覆盖有限的工具类型和攻击分布，现有结果不能替代更多真实业务流程中的长期效用与泛化评估。
